\documentclass[10pt,a4paper]{amsart}
\usepackage{amsmath,amssymb,amsthm,graphicx}
% Retained preamble for the Stacks Project chapter morphisms.tex (Morphisms of Schemes).
% The chapter file pinned-original-source.tex carries only \input{preamble} on its first line, so no
% preamble text lives inside the chapter itself.  The Stacks Project uses one single shared preamble for
% every chapter; the snapshot stages it verbatim in sources/stacks-chow/preamble.tex.  The block below is
% copied mechanically, line for line, from that file's lines 14-18 (the Xy-pic setup) and lines 154-217
% (the native theorem declarations together with the mathematical macros).  Nothing was added, removed or
% redefined.  The wrapper class and the amsmath/amssymb/amsthm/graphicx base packages are supplied by the
% assembler, so the class/\externaldocument part of the shared preamble is not repeated here.
% For commutative diagrams we use Xy-pic
\usepackage[all]{xy}

% We use 2cell for 2-commutative diagrams.
\xyoption{2cell}
% Theorem environments.
%
\theoremstyle{plain}
\newtheorem{theorem}[subsection]{Theorem}
\newtheorem{proposition}[subsection]{Proposition}
\newtheorem{lemma}[subsection]{Lemma}

\theoremstyle{definition}
\newtheorem{definition}[subsection]{Definition}
\newtheorem{example}[subsection]{Example}
\newtheorem{exercise}[subsection]{Exercise}
\newtheorem{situation}[subsection]{Situation}

\theoremstyle{remark}
\newtheorem{remark}[subsection]{Remark}
\newtheorem{remarks}[subsection]{Remarks}

\numberwithin{equation}{subsection}

% Macros
%
\def\lim{\mathop{\mathrm{lim}}\nolimits}
\def\colim{\mathop{\mathrm{colim}}\nolimits}
\def\Spec{\mathop{\mathrm{Spec}}}
\def\Hom{\mathop{\mathrm{Hom}}\nolimits}
\def\Ext{\mathop{\mathrm{Ext}}\nolimits}
\def\SheafHom{\mathop{\mathcal{H}\!\mathit{om}}\nolimits}
\def\SheafExt{\mathop{\mathcal{E}\!\mathit{xt}}\nolimits}
\def\Sch{\mathit{Sch}}
\def\Mor{\mathop{\mathrm{Mor}}\nolimits}
\def\Ob{\mathop{\mathrm{Ob}}\nolimits}
\def\Sh{\mathop{\mathit{Sh}}\nolimits}
\def\NL{\mathop{N\!L}\nolimits}
\def\CH{\mathop{\mathrm{CH}}\nolimits}
\def\proetale{{pro\text{-}\acute{e}tale}}
\def\etale{{\acute{e}tale}}
\def\QCoh{\mathit{QCoh}}
\def\Ker{\mathop{\mathrm{Ker}}}
\def\Im{\mathop{\mathrm{Im}}}
\def\Coker{\mathop{\mathrm{Coker}}}
\def\Coim{\mathop{\mathrm{Coim}}}

% Boxtimes
%
\DeclareMathSymbol{\boxtimes}{\mathbin}{AMSa}{"02}

%
% Macros for moduli stacks/spaces
%
\def\QCohstack{\mathcal{QC}\!\mathit{oh}}
\def\Cohstack{\mathcal{C}\!\mathit{oh}}
\def\Spacesstack{\mathcal{S}\!\mathit{paces}}
\def\Quotfunctor{\mathrm{Quot}}
\def\Hilbfunctor{\mathrm{Hilb}}
\def\Curvesstack{\mathcal{C}\!\mathit{urves}}
\def\Polarizedstack{\mathcal{P}\!\mathit{olarized}}
\def\Complexesstack{\mathcal{C}\!\mathit{omplexes}}
% \Pic is the operator that assigns to X its picard group, usage \Pic(X)
% \Picardstack_{X/B} denotes the Picard stack of X over B
% \Picardfunctor_{X/B} denotes the Picard functor of X over B
\def\Pic{\mathop{\mathrm{Pic}}\nolimits}
\def\Picardstack{\mathcal{P}\!\mathit{ic}}
\def\Picardfunctor{\mathrm{Pic}}
\def\Deformationcategory{\mathcal{D}\!\mathit{ef}}

% --- appended 2026-10-12 (Chief): the project own definitions of the reference/slogan/history
% display environments, copied mechanically from the shared Stacks preamble lines 8-12
% (\usepackage{verbatim} plus the three \newenvironment lines, which the chapter preamble keeps
% verbatim from the project preamble). Nothing added, removed or redefined: these environments
% carry no mathematical content and the project itself typesets them as comments. Without these
% three lines a counted statement whose own text uses them fails with Environment slogan undefined.
\usepackage{verbatim}
\newenvironment{reference}{\comment}{\endcomment}
\newenvironment{slogan}{\comment}{\endcomment}
\newenvironment{history}{\comment}{\endcomment}
\usepackage{cleveref}
\crefname{originallemma}{original lemma}{original lemmas}
\Crefname{originallemma}{Original lemma}{Original lemmas}
\makeatletter
\newlabel{algebra-lemma-help-with-powers}{{0545}{source}{}{sourcepointer-1}{}}
\newlabel{algebra-lemma-help-with-powers@cref}{{[originallemma][0][]0545}{[1][0][]source}{}{}{}}
\makeatother
\crefname{originallemma}{original lemma}{original lemmas}
\Crefname{originallemma}{Original lemma}{Original lemmas}
\makeatletter
\newlabel{algebra-lemma-image-dense-generic-points}{{00FL}{source}{}{sourcepointer-2}{}}
\newlabel{algebra-lemma-image-dense-generic-points@cref}{{[originallemma][0][]00FL}{[1][0][]source}{}{}{}}
\makeatother
\crefname{originallemma}{original lemma}{original lemmas}
\Crefname{originallemma}{Original lemma}{Original lemmas}
\makeatletter
\newlabel{algebra-lemma-p-ring-map}{{00I5}{source}{}{sourcepointer-3}{}}
\newlabel{algebra-lemma-p-ring-map@cref}{{[originallemma][0][]00I5}{[1][0][]source}{}{}{}}
\makeatother
\crefname{originallemma}{original lemma}{original lemmas}
\Crefname{originallemma}{Original lemma}{Original lemmas}
\makeatletter
\newlabel{lemma-pth-power-and-multiple}{{0CNC}{source}{}{sourcepointer-4}{}}
\newlabel{lemma-pth-power-and-multiple@cref}{{[originallemma][0][]0CNC}{[1][0][]source}{}{}{}}
\makeatother
\title{A ring map which becomes an isomorphism after inverting p is a universal homeomorphism exactly when its kernel is locally nilpotent and every element is bridged by a p-power multiple and a p-power}
\author{The Stacks Project Authors}
\date{stacks-project:morphisms.tex:local-snapshot-20261012; GNU FDL 1.2 or later}
\begin{document}
\maketitle
\noindent\textit{Editorial note.} All delivered mathematical text is the author's own, copied from the pinned Stacks Project chapter file pinned-original-source.tex; no statement or proof step was rewritten, reordered or removed.\par\medskip
\noindent\textit{Editorial note.} Following the multi-claim convention for this source, source\_url is the source-level URL https://stacks.math.columbia.edu/ and the source label lemma-universal-homeo-iso-if-invert-p is carried in claim\_key and here rather than in a per-claim tag URL.\par\medskip
\noindent\textit{Editorial note.} The preamble is the chapter's shared Stacks Project preamble, retained as sources/stacks-morphisms/retained-preamble.tex: a line-for-line excerpt of the shared Stacks preamble staged with this snapshot. No macro was invented or redefined, and this card adds no preamble text of its own.\par\medskip
\noindent\textit{Editorial note.} Statement span 11825-11837 and proof span 11839-11870 come from the halfA pool card operations/pool-cards/stacks-morphisms-halfA.json and were re-read in the pinned chapter file before this card was written.\par\medskip
\noindent\textit{Editorial note.} The statement's parenthesis giving the example of a nilpotent p is the author's and is preserved.\par\medskip
\noindent\textit{Editorial note.} The proof names the candidate subalgebra B' by an explicit set-builder expression; the display is preserved unchanged.\par\medskip
\noindent\textit{Editorial note.} 4 cross-references occur inside the counted spans; they are listed in external\_references as navigation pointers only.\par\medskip
\noindent\textit{Allowed standard prerequisites (editorial).}\begin{itemize}
\item The algebra criterion that a ring map whose kernel is locally nilpotent and which has the stated bridging property gives a universal homeomorphism. Cited at proof lines 11839-11840. Gives the easier implication from the algebraic condition to the geometric one.
\item If a morphism of spectra is a universal homeomorphism then the kernel of the corresponding ring map consists of nilpotent elements, cited from the algebra chapter. Cited at proof lines 11841-11843. Allows the base ring to be replaced by the image, so that the map is injective.
\item The set of elements bridged by a p-power multiple and a p-power is a subalgebra, cited from the algebra chapter. Cited at proof lines 11845-11849. Produces the subalgebra whose equality with the target is what has to be shown.
\item The two-step criterion producing an element outside a subalgebra with either two consecutive powers inside or a prime multiple and a prime power inside, cited from the same chapter. Cited at proof lines 11851-11854. Supplies the element used to derive the contradiction that closes the proof.
\end{itemize}\medskip
\ifdefined\hypertarget\hypertarget{sourcepointer-1}{}\fi
\noindent\textit{Source reference.} Original lemma 0545: focused-seeds/AGP/sources/stacks-snapshot-20261011/tags.txt, source lines 6628--6628. Same-chapter reference used in the retained author text at source line(s) 11844. It is supplied only as a navigation pointer so that the retained text compiles; the counted item delivers the author's own statement and the author's complete proof, and no part of the referenced item is copied in or substituted for the counted argument.\par\medskip
\ifdefined\hypertarget\hypertarget{sourcepointer-2}{}\fi
\noindent\textit{Source reference.} Original lemma 00FL: focused-seeds/AGP/sources/stacks-snapshot-20261011/tags.txt, source lines 575--575. Same-chapter reference used in the retained author text at source line(s) 1301, 11842, 11946, 12203, 14779. It is supplied only as a navigation pointer so that the retained text compiles; the counted item delivers the author's own statement and the author's complete proof, and no part of the referenced item is copied in or substituted for the counted argument.\par\medskip
\ifdefined\hypertarget\hypertarget{sourcepointer-3}{}\fi
\noindent\textit{Source reference.} Original lemma 00I5: focused-seeds/AGP/sources/stacks-snapshot-20261011/tags.txt, source lines 667--667. Same-chapter reference used in the retained author text at source line(s) 11820, 11840, 11973. It is supplied only as a navigation pointer so that the retained text compiles; the counted item delivers the author's own statement and the author's complete proof, and no part of the referenced item is copied in or substituted for the counted argument.\par\medskip
\ifdefined\hypertarget\hypertarget{sourcepointer-4}{}\fi
\noindent\textit{Source reference.} Original lemma 0CNC: focused-seeds/AGP/sources/stacks-snapshot-20261011/tags.txt, source lines 15887--15887. Same-chapter reference used in the retained author text at source line(s) 11807, 11851, 12205. It is supplied only as a navigation pointer so that the retained text compiles; the counted item delivers the author's own statement and the author's complete proof, and no part of the referenced item is copied in or substituted for the counted argument.\par\medskip
\begin{lemma}
\label{lemma-universal-homeo-iso-if-invert-p}
Let $p$ be a prime number. Let $A \to B$ be a ring map
which induces an isomorphism $A[1/p] \to B[1/p]$
(for example if $p$ is nilpotent in $A$).
The following are equivalent
\begin{enumerate}
\item $\Spec(B) \to \Spec(A)$ is a universal homeomorphism, and
\item the kernel of $A \to B$ is a locally nilpotent ideal and
for every $b \in B$ there exists a $p$-power $q$ with $qb$ and $b^q$
in the image of $A \to B$.
\end{enumerate}
\end{lemma}

\begin{proof}
If (2) holds, then (1) holds by Algebra, Lemma \ref{algebra-lemma-p-ring-map}.
Assume (1). Then the kernel of $A \to B$ consists of nilpotent
elements by Algebra, Lemma \ref{algebra-lemma-image-dense-generic-points}.
Thus we may replace $A$ by the image of $A \to B$ and assume that
$A \subset B$. By Algebra, Lemma \ref{algebra-lemma-help-with-powers}
the set
$$
B' = \{b \in B \mid p^nb, b^{p^n} \in A\text{ for some }n \geq 0\}
$$
is an $A$-subalgebra of $B$ (being closed under products is trivial).
We have to show $B' = B$. If not, then
according to Lemma \ref{lemma-pth-power-and-multiple}
there exists a $b \in B$, $b \not \in B'$ with either
$b^2, b^3 \in B'$ or there exists a prime number $\ell$
with $\ell b, b^\ell \in B'$.
We will show both cases lead to a contradiction, thereby proving the lemma.

\medskip\noindent
Since $A[1/p] = B[1/p]$ we can choose a $p$-power $q$ such that $qb \in A$.

\medskip\noindent
If $b^2, b^3 \in B'$ then also $b^q \in B'$. By definition of $B'$
we find that $(b^q)^{q'} \in A$ for some $p$-power $q'$.
Then $qq'b, b^{qq'} \in A$ whence $b \in B'$ which is a contradiction.

\medskip\noindent
Assume now there exists a prime number $\ell$ with $\ell b, b^\ell \in B'$.
If $\ell \not = p$ then $\ell b \in B'$ and $qb \in A \subset B'$ imply
$b \in B'$ a contradiction. Thus $\ell = p$ and $b^p \in B'$
and we get a contradiction exactly as before.
\end{proof}

\end{document}
